\section{Validated Benchmark and Experimental Protocol}
\label{sec:experimental-protocol}

\subsection{Benchmark Provenance and Numerical Gates}

The primary plant uses Julia 1.11.9, PowerDynamics.jl 5.0.0, and NetworkDynamics.jl 1.3.0 \cite{plietzsch2022powerdynamics}. It contains 39 buses, 46 branches, 19 loads, 10 sixth-order Sauer--Pai machines, nine AVR Type-I controllers, and nine TGOV1 governors. The descriptor realization has 192 coordinates: 114 differential and 78 algebraic. These parameters are package-defined tutorial data. We therefore classify the plant as a numerically validated pedagogical benchmark rather than an authoritative representation of an external grid.

The validation ladder prevents estimator performance from hiding an upstream modeling error. A same-plant static translation was audited after correcting tap orientation, shunts, line charging, and full-$Y_{\mathrm{bus}}$ injection accounting. Final maximum discrepancies were $1.14\times10^{-13}$~pu for a $Y_{\mathrm{bus}}$ element, $1.73\times10^{-7}$~pu for voltage magnitude, $1.40\times10^{-5}$ degrees for angle, and $1.71\times10^{-3}$~MVA for terminal flow. The controlled PMU operator then achieved maximum complex-voltage and analytic/backend terminal-current discrepancies of $2.71\times10^{-7}$ and $1.43\times10^{-14}$~pu, respectively.

Dynamic validation checks the equilibrium, local modes, nonlinear trajectories, and linearization. The DAE equilibrium residual is $4.66\times10^{-13}$. After removing the reference mode, the least-damped physical pair is $-0.09818\pm\mathrm{j}0.42948$~s$^{-1}$. Nonlinear Rodas5P simulations remain finite for the audited perturbations. Amplitude sweeps satisfy
\begin{equation}
 \|\bm y_{\mathrm{NL}}(\epsilon)-\bm y_{\mathrm{LIN}}(\epsilon)\|_2\propto\epsilon^p, \label{eq:linearization-order}
\end{equation}
with fitted slopes 1.9820, 1.9987, and 2.0003 for representative state, governor-reference, and load perturbations, consistent with first-order linearization error.

\begin{table*}[t]
\caption{Frozen experimental contracts. A complete nonlinear trajectory or rebuilt plant is the independent unit.}
\label{tab:experimental-contracts}
\centering
\footnotesize
\begin{tabularx}{\textwidth}{lXllll}
\toprule
Study & Purpose and data & Split & Units & Frames & Model use \\
\midrule
Structural audit & Local VI-only target residual for 31 hidden buses & preregistered & 31 buses & horizons 0--180 & frozen equilibrium \\
Nominal reconstruction & Governor, load, and combined nonlinear perturbations & 20 DEV / 100 TEST & trajectories & 91 at 30 Hz & frozen B0--B3 \\
Nominal calibration & New nonlinear nominal trajectories & 30 CAL & trajectories & 91 at 30 Hz & fit U1/U2 and discrepancy only \\
Physical mismatch & Seven families, seven scales, 20 seeds & frozen TEST & 980 rebuilt plants & 91 at 30 Hz & frozen B1/B2/U2 \\
Breakpoint refinement & New seeds in transition bands & frozen TEST & 175 rebuilt plants & 91 at 30 Hz & frozen B2 \\
Causal recentering & M1/M2/M6/M7 with disjoint new seeds & 160 DEV / 320 TEST & rebuilt plants & 91 at 30 Hz & select on DEV; score TEST once \\
\bottomrule
\end{tabularx}
\end{table*}

\subsection{Nominal Reconstruction and Calibration}

The nominal dataset contains 20 development and 100 test trajectories, each sampled at 30~Hz for 3~s. Perturbations are normal-operation changes to governor references, loads, or their combination. The equilibrium, linear model, $\bm P_0=10^{-2}\bm I$, $\bm Q=10^{-6}\bm I$, and $\bm R=10^{-6}\bm I$ are frozen before TEST. B0, B1, B2, and all B3 lags are evaluated on the same trajectories. A separate 30-trajectory calibration split fits U1/U2 temperatures and low-rank discrepancy variants; TEST is not used for selection.

Hidden truth is used only for post-hoc metrics. The eight observed buses are excluded from every primary reconstruction score. The phasor prediction is formed as $\bm V_0+\Delta\bm V$, and angle error is wrapped using $\operatorname{atan2}(\sin\Delta\vartheta,\cos\Delta\vartheta)$. These rules correct an earlier exploratory implementation that incorrectly interpreted $\angle\Delta V$ as the absolute angle; no superseded result from that implementation is retained.

\subsection{Physical Mismatch Campaign}

For each mismatch case, a fresh PowerDynamics plant is rebuilt from copied source tables, followed by power flow, dynamic initialization, and nonlinear simulation. The original package CSVs remain hash-identical. The seven families are network $R/X/B$ (M1), machine parameters (M2), governor parameters (M3), AVR parameters (M4), ZIP/load model (M5), operating-point redispatch (M6), and coupled mismatch (M7). Stress scale is $m\in\{0,0.25,0.5,0.75,1.0,1.25,1.5\}$. These are engineering stress-test ranges, not probability priors.

All 980 main-grid cases and all 175 refinement cases passed power-flow, initialization, and time-domain acceptance. B1, B2, U0, U2, and the adequacy score are frozen from nominal data. Bootstrap intervals use seeds as independent units. Breakpoints in this paper are not inferred from the narrative report: the generated plots use the frozen per-case table directly because a prior refined breakpoint mixed aggregation definitions.

\subsection{Oracle and Causal Adaptation Contracts}

The oracle subset compares the frozen estimator with true-equilibrium recentering and true local relinearization on the same mutated trajectories. Oracle quantities are evaluation-only. The causal recentering study uses new DEV seeds 101--110 and TEST seeds 201--220, disjoint from all earlier datasets. Its window, cadence, regularization, and activation threshold are selected on DEV and frozen before TEST. No event labels, hidden state, or mismatch family enter the R2-A estimate. Rejected or nonconvergent cases would remain in the denominator; none occurred in the reported physical mismatch grids.

All numerical tables and empirical figures are regenerated from the hashed snapshot under \texttt{generated/frozen}. The build script asserts 100 nominal test trajectories, 980 main-grid B2 cases, and the expected adaptation closure range before producing manuscript assets.
